\subsubsection{Mixer Hamiltonian and Mixer Operator}\label{sec:mixer-hamiltonian-and-mixer-operator}

In the previous section we saw the \textbf{Cost Hamiltonian}, whose purpose is to \textbf{encode the optimization problem}. Now we introduce the Mixer Hamiltonian. Its purpose is \textbf{not} to optimize the cost function, but to \textbf{continuously move the quantum state through the search space}, preventing the algorithm from getting trapped in a fixed configuration.

\highspace
The Mixer Hamiltonian is fundamental. Without it, the Cost Operator would repeatedly evolve the state according to the same energy landscape. The problem is that \textbf{nothing would encourage the algorithm to explore different candidate solutions}. Instead, the state could become concentrated around a limited region of the search space. Instead, the \definition{Mixer Hamiltonian} solves exactly this problem. Its purpose is \textbf{to explore the space of possible quantum states (candidate solutions)}.

\highspace
\begin{flushleft}
    \textcolor{Green3}{\faIcon{balance-scale} \textbf{Cost vs Mixer}}
\end{flushleft}
The role of the Cost and Mixer Hamiltonians are complementary. Let's summarize their differences in the following table:

\begin{table}[!htp]
    \centering
    \begin{tabular}{@{} l l @{}}
        \toprule
        \textbf{Cost Hamiltonian} & \textbf{Mixer Hamiltonian} \\
        \midrule
        Encodes the optimization problem    & Explores the search space             \\
        Depends on the problem              & Usually fixed                         \\
        Built from QUBO/Ising               & Built from Pauli-$X$ operators        \\
        Tries to favor good solutions       & Creates transitions between solutions \\
        \bottomrule
    \end{tabular}
\end{table}

\highspace
\begin{flushleft}
    \textcolor{Green3}{\faIcon{book} \textbf{Mixer Hamiltonian}}
\end{flushleft}
The standard QAOA mixer is:
\begin{equation}
    H_M = \sum_j \sigma_{X}^{j}
\end{equation}
Where each $\sigma_{X}^{j}$ \hl{applies a Pauli-$X$ operator only to the $j$-th qubit, while the identity acts on all the others:}
\begin{equation*}
    \sigma_{X}^{j} = I \otimes I \otimes \dots \otimes \underbrace{\sigma_{X}}_{j\text{-th}} \otimes \dots \otimes I
\end{equation*}
For example, with three qubits:
\begin{equation*}
    \begin{aligned}
        \sigma_{X}^{1} &= \sigma_{X} \otimes I \otimes I \\
        \sigma_{X}^{2} &= I \otimes \sigma_{X} \otimes I \\
        \sigma_{X}^{3} &= I \otimes I \otimes \sigma_{X}
    \end{aligned}
\end{equation*}
Therefore, the \textbf{Mixer Hamiltonian} is simply a sum of one Pauli-$X$ operator for every qubit:
\begin{equation*}
    H_M = X_1 + X_2 + \dots + X_n
\end{equation*}
\textcolor{Green3}{\faIcon{question-circle} \textbf{Why Pauli-$X$?}} The Pauli-$X$ operator is a quantum gate that flips the state of a qubit. In the context of QAOA, suppose our current candidate solution is:
\begin{equation*}
    01010
\end{equation*}
Applying an $X$ operation on one qubit changes one variable in the candidate solution. For example, applying $X$ on the first qubit gives:
\begin{equation*}
    11010
\end{equation*}
The Mixer Hamiltonian therefore \textbf{creates transitions between different bit strings}. It allows the algorithm to leave the current candidate solution and explore nearby ones (i.e., those that differ by a single bit flip). This is why the mixer is associated with \textbf{exploration}. Unlike the Cost Hamiltonian, which only assigns energies, the \textbf{Mixer Hamiltonian actively mixes amplitudes between computational basis states}.

\highspace
\begin{flushleft}
    \textcolor{Green3}{\faIcon{book} \textbf{The Mixer Operator}}
\end{flushleft}
Exactly as we did for the Cost Hamiltonian, we can define the \definition{Mixer Operator} as the unitary evolution generated by the Mixer Hamiltonian:
\begin{equation}\label{eq:mixer-operator}
    U_M\left(\beta_i\right) = e^{-i \beta_i H_M} \quad \text{with } \beta_i \in \mathbb{R}
\end{equation}
\textcolor{Green3}{\faIcon{question-circle} \textbf{Why is the Mixer Hamiltonian fixed?}} Unlike the Cost Hamiltonian, which depends on the optimization problem, the \textbf{Mixer Hamiltonian is usually the same for every problem}. Because exploration is a generic operation. Regardless of whether we solve a Max-Cut problem, a Traveling Salesman Problem, or any other combinatorial optimization problem, we still \textbf{want to move through the space of candidate solutions}. Therefore, the \hl{same exploration mechanism works for many optimization problems.}

\highspace
\textcolor{Green3}{\faIcon{tools} \textbf{How is it implemented?}} Unlike the Cost Hamiltonian, the Mixer Hamiltonian is already extremely simple to implement. Since it is a sum of Pauli-$X$ operators:
\begin{equation*}
    H_M = \sum_j \sigma_{X}^{j}
\end{equation*}
And we already know how to implement the unitary evolution of a Pauli-$X$ operator:
\begin{equation*}
    R_{x}\left(\theta\right) = e^{-i \frac{\theta}{2} \sigma_{X}}
\end{equation*}
Each $\sigma_{X}^{j}$ term corresponds to an $R_{x}$ rotation on the $j$-th qubit. Therefore, the Mixer Operator is \hl{implemented as a layer of parallel $R_{x}\left(\beta_i\right)$ gates applying the same rotation angle to every qubit.} With ``\emph{parallel}'' we mean that all the $R_{x}$ gates can be \textbf{applied simultaneously}, since they act on different qubits. This parallelism means also that $\beta_i$ is a \textbf{single parameter for the entire layer}, not one per qubit. The Mixer Operator is therefore implemented as a single layer of $R_{x}$ gates, all with the same rotation angle $\beta_i$. The optimizer will search for a vector of angles $\vec{\beta} = \left(\beta_1, \beta_2, \dots, \beta_p\right)$ for the $p$ layers of the QAOA ansatz.